\documentclass[conference]{IEEEtran}

\usepackage{cite}
\usepackage{amsmath,amssymb,amsfonts}
\usepackage{graphicx}
\usepackage{textcomp}
\usepackage{xcolor}
\usepackage{booktabs}
\usepackage{url}

\begin{document}

% =========================================================
% TITLE AND AUTHOR
% =========================================================

\title{Reliable Machine-Failure Prediction Under Imbalanced and Corrupted Sensor Data}

\author{
\IEEEauthorblockN{MAA AACAAS MUHAMATH}
\IEEEauthorblockA{
Department of Electrical Engineering\\
University of Moratuwa\\
Moratuwa, Sri Lanka
}
}

\maketitle

% =========================================================
% ABSTRACT
% =========================================================

\begin{abstract}

Machine-learning systems for predictive maintenance are commonly evaluated under clean test conditions, although real industrial systems may experience noisy sensors, missing measurements, class imbalance, and limited labelled failure data.

This study investigates machine-failure prediction reliability under these conditions using the AI4I 2020 Predictive Maintenance Dataset \cite{uci_ai4i}. A Dummy Classifier, Logistic Regression, and Random Forest were evaluated during baseline modelling. An unweighted Random Forest containing 300 trees was selected as the final prediction system, with a validation-selected decision threshold of 0.20.

Under clean test conditions, the selected model achieved 97.27\% accuracy, 56.76\% precision, 82.35\% recall, 67.20\% F1-score, and 78.03\% Average Precision.

Robustness was evaluated under increasing sensor noise, missing sensor measurements, and reduced labelled training data. At 30\% sensor noise, recall decreased to 66.67\%, while 30\% missing sensor measurements reduced recall to 50.98\%. Training with only 20\% of the available labelled data resulted in approximately 66.67\% recall.

Model explainability was investigated using Random Forest feature importance and SHAP \cite{lundberg2017}. Torque represented approximately 32.45\% of total impurity-based feature importance, while rotational speed and tool wear were also influential.

The results demonstrate that high clean-data accuracy alone is insufficient for evaluating predictive-maintenance reliability. Class-imbalance-aware metrics, robustness testing, and explainability provide a more comprehensive assessment.

\end{abstract}

\begin{IEEEkeywords}
predictive maintenance, machine failure prediction, Random Forest, class imbalance, sensor noise, missing data, robust machine learning, SHAP, explainable artificial intelligence
\end{IEEEkeywords}

\noindent
\textbf{Project Repository:}
\url{https://github.com/aacaas5/reliable-machine-failure-prediction-under-imbalanced-and-corrupted-sensor-data}

% =========================================================
% I. INTRODUCTION
% =========================================================

\section{Introduction}

Predictive maintenance aims to identify equipment degradation or impending failure before breakdown occurs. Machine-learning techniques can support this objective by analysing multiple operating measurements simultaneously to identify patterns associated with equipment failure.

However, predictive-maintenance datasets frequently exhibit strong class imbalance because normal machine operation occurs substantially more frequently than failure events.

Under these conditions, high accuracy can provide an overly optimistic interpretation of model performance. Precision--recall based metrics are particularly useful for evaluating minority-class performance in highly imbalanced binary classification problems \cite{saito2015}.

Real industrial deployments also introduce additional challenges. Sensors may become noisy, measurements may be unavailable, communication interruptions may occur, and labelled failure examples may be limited.

Therefore, a machine-learning model that performs strongly under clean benchmark conditions may not necessarily remain reliable during deployment.

This study investigates machine-failure prediction as a reliability problem rather than focusing only on clean-data classification performance.

The main research questions are:

\begin{enumerate}

    \item Can the selected machine-learning model detect failures effectively under clean-data conditions?

    \item How does failure-detection reliability change when sensor measurements become noisy?

    \item How does missing sensor information affect machine-failure prediction?

    \item What happens when the amount of labelled training data is reduced?

    \item Which features have the greatest influence on model predictions?

    \item Why are some actual failures detected confidently while others are missed?

\end{enumerate}

The main contribution of this project is a reliability-focused evaluation framework combining clean-data performance, robustness testing, class-imbalance-aware metrics, and model explainability.

% =========================================================
% II. DATASET
% =========================================================

\section{Dataset}

The AI4I 2020 Predictive Maintenance Dataset from the UCI Machine Learning Repository was used in this study \cite{uci_ai4i}.

The dataset contains 10,000 observations and was generated to reflect industrial predictive-maintenance conditions.

Six original input variables were used:

\begin{itemize}

    \item Machine type
    \item Air temperature
    \item Process temperature
    \item Rotational speed
    \item Torque
    \item Tool wear

\end{itemize}

The binary target variable represents machine failure:

\[
y =
\begin{cases}
0, & \text{no machine failure},\\
1, & \text{machine failure}.
\end{cases}
\]

The dataset was divided into:

\begin{itemize}

    \item 7000 training observations,
    \item 1500 validation observations,
    \item 1500 test observations.

\end{itemize}

The final test set contained 51 actual machine failures.

The target distribution was strongly imbalanced, making metrics such as recall, F1-score, and Average Precision particularly important.

% =========================================================
% III. METHODOLOGY
% =========================================================

\section{Methodology}

\subsection{Preprocessing}

A reusable preprocessing pipeline was used throughout the project.

Numerical variables were retained as numerical values for the tree-based model, while machine type was encoded into indicator variables.

After preprocessing, the Random Forest used eight transformed input features.

All preprocessing transformations were fitted using the training data only.

Validation data, clean test data, and corrupted test data were processed using the already-fitted preprocessing pipeline.

This approach prevented information from validation or test observations from influencing preprocessing parameters and reduced the risk of data leakage.

\subsection{Baseline Models}

Three baseline approaches were evaluated:

\begin{itemize}

    \item Dummy Classifier,
    \item Logistic Regression,
    \item Random Forest.

\end{itemize}

Random Forests combine multiple decision trees and are capable of modelling nonlinear relationships and interactions between variables \cite{breiman2001}.

The final selected model was an unweighted Random Forest configured with 300 trees.

\begin{verbatim}
RandomForestClassifier(
    n_estimators=300,
    random_state=42,
    n_jobs=-1
)
\end{verbatim}

The random state was fixed at 42 to improve reproducibility.

\subsection{Decision Threshold}

The Random Forest outputs an estimated probability of machine failure.

Instead of using the default probability threshold of 0.50, a decision threshold was selected using validation data.

The selected threshold was:

\[
T = 0.20.
\]

The prediction rule was therefore:

\[
\hat{y} =
\begin{cases}
1, & P(y=1) \geq 0.20,\\
0, & P(y=1) < 0.20.
\end{cases}
\]

The threshold was fixed before final test evaluation.

The same threshold was subsequently used during all robustness experiments and explainability analyses.

Maintaining a fixed threshold allowed direct comparison between clean and corrupted operating conditions.

\subsection{Evaluation Metrics}

Performance was evaluated using:

\begin{itemize}

    \item Accuracy,
    \item Precision,
    \item Recall,
    \item F1-score,
    \item Average Precision,
    \item Confusion matrix.

\end{itemize}

Recall was calculated as:

\[
\mathrm{Recall}
=
\frac{TP}{TP+FN}
\]

where $TP$ represents true positives and $FN$ represents false negatives.

Precision was calculated as:

\[
\mathrm{Precision}
=
\frac{TP}{TP+FP}.
\]

The F1-score was:

\[
F_1 =
2
\frac{
\mathrm{Precision}
\times
\mathrm{Recall}
}{
\mathrm{Precision}
+
\mathrm{Recall}
}.
\]

Precision--recall based evaluation is particularly useful in highly imbalanced binary classification problems \cite{saito2015}.

% =========================================================
% IV. CLEAN BASELINE
% =========================================================

\section{Clean-Data Baseline Results}

The selected Random Forest was first evaluated on the untouched test set.

The clean-data performance is shown in Table~\ref{tab:baseline}.

\begin{table}[htbp]

\caption{Clean-Data Random Forest Performance}

\label{tab:baseline}

\centering

\begin{tabular}{lc}

\toprule

\textbf{Metric} & \textbf{Result}\\

\midrule

Accuracy & 97.27\%\\
Precision & 56.76\%\\
Recall & 82.35\%\\
F1-score & 67.20\%\\
Average Precision & 78.03\%\\
Decision Threshold & 0.20\\

\bottomrule

\end{tabular}

\end{table}

The confusion matrix was:

\[
\begin{bmatrix}
1417 & 32\\
9 & 42
\end{bmatrix}.
\]

Among the 51 actual failures:

\begin{itemize}

    \item 42 failures were correctly detected,
    \item 9 failures were missed.

\end{itemize}

Although overall accuracy reached 97.27\%, the highly imbalanced class distribution means that accuracy alone provides an incomplete description of system reliability.

Recall and Average Precision provide more meaningful information about actual failure detection.

% =========================================================
% V. ROBUSTNESS EXPERIMENTS
% =========================================================

\section{Robustness Experiments}

Three robustness conditions were investigated:

\begin{enumerate}

    \item Sensor measurement noise
    \item Missing sensor measurements
    \item Limited labelled training data

\end{enumerate}

\subsection{Sensor Noise}

Random noise was introduced into numerical sensor measurements at severity levels of:

\[
0\%,\quad
5\%,\quad
10\%,\quad
20\%,\quad
30\%.
\]

Noise magnitude was scaled relative to the standard deviation of each numerical variable.

This prevented features operating at different numerical scales from receiving inappropriate absolute perturbations.

Throughout the experiment, the same:

\begin{itemize}

    \item test observations,
    \item true labels,
    \item fitted preprocessing pipeline,
    \item Random Forest model,
    \item decision threshold

\end{itemize}

were retained.

At 30\% sensor noise, performance deteriorated substantially.

\begin{table}[htbp]

\caption{Clean Versus 30\% Sensor Noise}

\centering

\begin{tabular}{lcc}

\toprule

\textbf{Metric} &
\textbf{Clean} &
\textbf{30\% Noise}\\

\midrule

Accuracy & 97.27\% & 93.60\%\\
Precision & 56.76\% & 30.09\%\\
Recall & 82.35\% & 66.67\%\\
F1-score & 67.20\% & 41.46\%\\
Average Precision & 78.03\% & 54.17\%\\

\bottomrule

\end{tabular}

\end{table}

Recall decreased from 82.35\% to 66.67\%.

Precision also decreased substantially, indicating an increase in false failure alarms.

\subsection{Missing Sensor Measurements}

The second robustness experiment simulated unavailable numerical sensor measurements.

Random sensor cells were replaced with missing values.

A missing-data level of 10\%, for example, represented approximately 10\% of individual numerical sensor measurements being unavailable.

Median values calculated from the training set were used for imputation.

Importantly, test-set statistics were not used to calculate the imputation values.

At 30\% missing sensor measurements:

\begin{itemize}

    \item recall decreased from 82.35\% to 50.98\%,
    \item accuracy remained approximately 96.87\%.

\end{itemize}

This result demonstrates an important limitation of accuracy in highly imbalanced machine-failure classification.

Even though the model missed approximately half of the real failures, overall accuracy remained high because healthy machines represented the majority of observations.

\subsection{Limited Training Data}

The third robustness experiment investigated the effect of reducing the quantity of labelled training data.

The Random Forest was retrained using stratified subsets corresponding to:

\[
20\%,\quad
40\%,\quad
60\%,\quad
80\%,\quad
100\%
\]

of the original training set.

The Random Forest configuration, clean test set, and decision threshold remained unchanged.

Average Precision improved as more labelled data became available.

\begin{table}[htbp]

\caption{Average Precision Versus Training Data}

\centering

\begin{tabular}{cc}

\toprule

\textbf{Training Fraction} &
\textbf{Average Precision}\\

\midrule

20\% & 65.46\%\\
40\% & 69.78\%\\
60\% & 74.84\%\\
80\% & 77.81\%\\
100\% & 78.03\%\\

\bottomrule

\end{tabular}

\end{table}

At the 20\% training-data condition, recall was approximately 66.67\%.

These results indicate that increasing the quantity of labelled training data improved the model's ability to rank and identify machine failures.

% =========================================================
% VI. ROBUSTNESS SUMMARY
% =========================================================

\section{Robustness Summary}

Representative recall values are shown in Table~\ref{tab:summary}.

\begin{table}[htbp]

\caption{Recall Under Representative Conditions}

\label{tab:summary}

\centering

\begin{tabular}{lc}

\toprule

\textbf{Condition} &
\textbf{Recall}\\

\midrule

Clean data & 82.35\%\\
30\% sensor noise & 66.67\%\\
30\% missing measurements & 50.98\%\\
20\% training data & 66.67\%\\

\bottomrule

\end{tabular}

\end{table}

Among the severe conditions tested, missing sensor measurements produced the largest reduction in failure recall.

The results demonstrate that sensor quality, sensor availability, and labelled data availability all influence failure-detection reliability.

% =========================================================
% VII. EXPLAINABILITY
% =========================================================

\section{Model Explainability}

Model explainability was investigated using two complementary approaches:

\begin{enumerate}

    \item Random Forest impurity-based feature importance
    \item SHAP analysis

\end{enumerate}

\subsection{Random Forest Feature Importance}

The impurity-based feature-importance analysis identified torque as the most influential variable.

Torque represented approximately:

\[
32.45\%
\]

of the Random Forest's total impurity-based feature importance.

Rotational speed and tool wear also contributed substantially.

Temperature variables showed moderate importance, while encoded machine-type variables contributed relatively little.

Feature importance describes the model's learned predictive behaviour and should not be interpreted as evidence of physical causation.

\subsection{SHAP Analysis}

SHAP provides additive feature-attribution explanations and can be used to analyse global and individual model behaviour \cite{lundberg2017}.

For the failure class:

\begin{itemize}

    \item positive SHAP values push the prediction toward failure,
    \item negative SHAP values push the prediction away from failure,
    \item larger absolute SHAP values indicate stronger influence.

\end{itemize}

Global SHAP analysis showed that rotational speed, torque, tool wear, and temperature measurements were among the most influential variables.

Machine-type features generally produced smaller SHAP effects.

% =========================================================
% VIII. INDIVIDUAL PREDICTIONS
% =========================================================

\section{Individual Prediction Analysis}

\subsection{Correctly Detected Failure}

One correctly detected failed machine received:

\[
P(\mathrm{failure}) = 0.94.
\]

This was substantially above the decision threshold of 0.20.

Its strongest SHAP contributions were:

\begin{table}[htbp]

\caption{SHAP Contributions for True-Positive Example}

\centering

\begin{tabular}{lcc}

\toprule

\textbf{Feature} &
\textbf{Value} &
\textbf{SHAP}\\

\midrule

Rotational speed & 2737 & +0.5036\\
Torque & 8.8 & +0.3625\\
Tool wear & 142 & +0.0087\\
Process temperature & 310.2 & +0.0075\\
Air temperature & 298.9 & +0.0021\\

\bottomrule

\end{tabular}

\end{table}

Rotational speed and torque strongly dominated this prediction.

\subsection{False-Negative Failure}

A selected false-negative machine belonged to the failure class but received:

\[
P(\mathrm{failure}) = 0.1367.
\]

Since:

\[
0.1367 < 0.20,
\]

the system classified the observation as no failure.

Its strongest SHAP contributions were:

\begin{table}[htbp]

\caption{SHAP Contributions for False-Negative Example}

\centering

\begin{tabular}{lcc}

\toprule

\textbf{Feature} &
\textbf{Value} &
\textbf{SHAP}\\

\midrule

Rotational speed & 1371 & +0.0910\\
Air temperature & 303.6 & +0.0906\\
Process temperature & 312.2 & -0.0684\\
Tool wear & 112 & -0.0206\\
Torque & 54.6 & +0.0143\\

\bottomrule

\end{tabular}

\end{table}

The false-negative case contained weaker and partly conflicting evidence.

Rotational speed and air temperature pushed the prediction toward failure, while process temperature and tool wear pushed it away from failure.

Torque also contributed substantially less than in the correctly detected failure.

% =========================================================
% IX. DISCUSSION
% =========================================================

\section{Discussion}

The experiments demonstrate that strong clean-data performance does not guarantee reliable behaviour under imperfect operating conditions.

The selected Random Forest achieved 82.35\% recall under clean conditions and successfully detected 42 of the 51 actual test failures.

However, recall decreased to 66.67\% under severe sensor noise and to 50.98\% under severe missing-data corruption.

The missing-data experiment is particularly important because overall accuracy remained approximately 96.87\%.

This demonstrates how an imbalanced target distribution can hide substantial deterioration in minority-class failure detection.

The limited-training-data experiment also showed that Average Precision improved as more labelled training examples became available.

Average Precision increased from 65.46\% when using only 20\% of the original training data to 78.03\% when using the complete training set.

Explainability analysis further demonstrated that numerical operating measurements played a greater role than machine-type variables.

The true-positive example was dominated by strong rotational-speed and torque contributions.

In comparison, the false-negative case contained weaker and conflicting evidence.

These observations illustrate why a model may perform strongly overall while still missing certain genuine failures.

% =========================================================
% X. LIMITATIONS AND FUTURE WORK
% =========================================================

\section{Limitations and Future Work}

Several limitations should be acknowledged.

First, this study evaluated a single benchmark dataset. The AI4I dataset is synthetic, and exact numerical performance may therefore not generalise directly to physical industrial environments.

Second, sensor noise and missing measurements were artificially simulated.

Real industrial sensor problems may involve:

\begin{itemize}

    \item persistent drift,
    \item calibration bias,
    \item correlated measurement errors,
    \item communication failures,
    \item complete sensor dropout.

\end{itemize}

Third, the reduced-training-data experiment used individual stratified subsets.

Because failure observations are rare, performance may vary depending on which minority-class examples are included.

Fourth, the detailed robustness evaluation focused primarily on one selected Random Forest system.

Future work could investigate:

\begin{itemize}

    \item real industrial predictive-maintenance datasets,
    \item repeated robustness experiments across multiple random seeds,
    \item persistent sensor drift,
    \item correlated sensor corruption,
    \item complete sensor dropout,
    \item alternative missing-value strategies,
    \item gradient-boosted tree models,
    \item neural networks,
    \item temporal models,
    \item probability calibration,
    \item uncertainty estimation,
    \item cost-sensitive learning,
    \item application-specific threshold optimisation.

\end{itemize}

An important future extension would be uncertainty-aware prediction so that the system can identify situations where insufficient sensor information makes a failure prediction unreliable.

% =========================================================
% XI. CONCLUSION
% =========================================================

\section{Conclusion}

This study investigated machine-failure prediction as a reliability problem rather than evaluating classification performance only under clean-data conditions.

The selected Random Forest achieved strong clean-data performance, including 82.35\% recall and 78.03\% Average Precision using a fixed decision threshold of 0.20.

However, controlled robustness experiments demonstrated that sensor corruption, missing measurements, and limited labelled training data reduced failure-detection reliability.

Missing sensor measurements produced the largest reduction in recall among the severe tested conditions.

Explainability analysis further demonstrated that model predictions depended on combinations of operating measurements rather than a single variable.

Correctly detected failures could exhibit strong learned feature patterns, whereas missed failures could contain weaker or conflicting model evidence.

The central conclusion is therefore that predictive-maintenance models should not be judged using clean-data accuracy alone.

A more reliable evaluation should combine:

\begin{itemize}

    \item class-imbalance-aware metrics,
    \item controlled robustness testing,
    \item consistent decision thresholds,
    \item analysis of missed failures,
    \item interpretable model behaviour.

\end{itemize}

% =========================================================
% REFERENCES
% =========================================================

\begin{thebibliography}{00}

\bibitem{uci_ai4i}
S. Matzka,
``AI4I 2020 Predictive Maintenance Dataset,''
UCI Machine Learning Repository, 2020.
doi: 10.24432/C5HS5C.

\bibitem{breiman2001}
L. Breiman,
``Random Forests,''
\textit{Machine Learning},
vol. 45, no. 1,
pp. 5--32,
2001.
doi: 10.1023/A:1010933404324.

\bibitem{lundberg2017}
S. M. Lundberg and S.-I. Lee,
``A Unified Approach to Interpreting Model Predictions,''
in \textit{Advances in Neural Information Processing Systems},
vol. 30,
2017.

\bibitem{saito2015}
T. Saito and M. Rehmsmeier,
``The Precision-Recall Plot Is More Informative than the ROC Plot When Evaluating Binary Classifiers on Imbalanced Datasets,''
\textit{PLOS ONE},
vol. 10,
no. 3,
Art. no. e0118432,
2015.
doi: 10.1371/journal.pone.0118432.

\end{thebibliography}

\end{document}
